% =====================================================================
%  Abstract: English
% =====================================================================
\chapter*{Abstract}
\addcontentsline{toc}{chapter}{Abstract}
\thispagestyle{plain}

This thesis describes the design, the implementation and the deployment of an end-to-end data engineering and machine learning platform built for the multi-tenant telematics fleet of \hostorg{}. The host organization operates three MySQL databases that fulfil complementary roles in the lifecycle of a telemetry signal: a raw-stream database in which the frames pushed by GPS/IoT devices through SIM-card connectivity are stored as hexadecimal payloads; a decoded archive whose tables follow the convention \texttt{arch\_xxxxx} and retain approximately the most recent six months of readable telemetry; and a master-data database whose tables follow the convention \texttt{user\_xxxxx} and carry the per-tenant customers, vehicles and drivers. Although the company already exposes mature dashboards and standard reports on top of these sources, the analytical value of the historical archive remains largely untapped: there is no behavioural segmentation, no anomaly score and no continuous risk indicator that would let the operations team anticipate operational risk rather than only observe past activity. Following the CRISP-DM methodology, the project introduces an analytical layer that explicitly addresses this gap. Selected slices of the operational MySQL data are exported as \texttt{.sql} dump files, converted by a dedicated MySQL-to-PostgreSQL bridge into a PostgreSQL-importable representation, and loaded into a three-layer warehouse (staging, warehouse, marts) hosted on Microsoft Azure. The choice of PostgreSQL~16 as the analytical DBMS is justified by its strong support for analytical workloads, its ACID compliance, its advanced indexing, its scalability to large datasets, its extensibility and its first-class compatibility with the Python data ecosystem and standard BI tools. Because the project is conducted by a two-student team of engineers working from separate workstations, the analytical environment is hosted on an Azure virtual machine accessed through SSH tunnelling, with an IP allowlist enforced at the Network Security Group level. The thesis covers business understanding, data understanding through PostgreSQL profiling, data preparation through an explicit cleaning rule engine, feature engineering of 35 behavioural indicators at the device-month grain, the training of a tenant-aware K-Means segmentation, the evaluation of an Isolation Forest anomaly indicator with explicit caveats, and finally the production deployment of the transparent risk profile through a Prefect flow scheduled by cron, served through a FastAPI service and a Streamlit dashboard.

\noindent\textbf{Keywords:} CRISP-DM, telematics, fleet analytics, advanced analytics, MySQL, PostgreSQL, data engineering, ETL, feature engineering, K-Means, Isolation Forest, Microsoft Azure, SSH tunnelling, monitoring.

\bigskip
\hrule
\bigskip

\cleardoublepage
